\section{Keeping the population and comparison fixed}
\label{sec:exec28-panel-delegates}

The preceding design comparison becomes actionable only when the selected software carries the same baseline, support
and event order. Ten further pinned files expose those decisions in \texttt{fastdid}, \texttt{nonabsdid} and \texttt{pampe}.
Their value is not that a function returns a number, but that its observable selection and transformation rules can be
compared with the banking question. This is a static source comparison with independent arithmetic, not execution of the
acquired packages.\citep{exec28_panel_delegates}

\subsection{The multiple-event anchor is explicit in the implementation}

In the selected \texttt{fastdid} commit, the event-specific weight routine separates a direct case before the competing
event, an imputation case when the target arrives first, and a difference-in-differences case when it arrives later.
For target-first imputation it assigns the deterministic weight to the target cohort at the pre-confounding base period,
then adds the eligible donor change. Thus this inspected branch agrees with the explicit Appendix A.4 anchor developed
earlier, rather than equation (11)'s inconsistent target-period entry in the selected v3. That resolves one concrete
implementation question; it does not resolve the different edition named by the repository or verify its general
multiple-event theorem.\citep{exec28_panel_delegates}

The routine also intersects donor cohort identities available at both periods. This matters independently of the anchor.
Suppose baseline donors A and B have effects 10 and 20, while outcome donors B and C have effects 25 and 60. Separately
normalizing available donors gives a change of $55/2$; the common donor B changes by only 5. Transporting $55/2$ would
mix a change in composition with a change in effects. The inspected code restricts both sides to the common identities
and warns and skips the cell if none remain. Its target-later branch matches already-arrived competing event dates and
requires other events to remain sufficiently far away. At equal first/earliest-competing event dates, no separating branch
supplies a nonzero weight. These are useful refusal semantics to retain in a bank-facing response.\citep{exec28_panel_delegates,exec28_panel_delegate_reference}

The aggregation routine transforms both effects and influence functions, including estimated cohort-share contributions
in the second-stage and final aggregation. That is stronger than merely multiplying reported standard errors by weights.
Nevertheless its covariance depends on downstream bootstrap and estimation delegates not executed here. Selected validation
rejects doubly robust estimation with an unbalanced panel, time-varying covariates with an unbalanced panel, and additional
clustering or uniform bands without bootstrap. It may drop units with missing competing-event dates, incomplete panels or
first-period exposure. The first-stage driver can skip failed or very small cells. An adapter must report those exclusions
and the resulting target population; it must not quietly substitute zeros for refused cells.\citep{exec28_panel_delegates}

Two additional checks remain important. Calendar normalization requires uniformly spaced periods and shifts every event
date consistently. The dynamic balance selector checks cohort extrema, but that alone is not proof that every interior
event-time cell survived earlier estimation failures. Full membership checks are therefore required for a promised common
window. The two-stage control routine uses both \texttt{double\_control\_option} and \texttt{control\_option} in different
branches; their combinations need an explicit fixture matrix before the options are exposed as interchangeable concepts.
The proposed reuse is a pinned reference integration with visible support and exclusions, conditional on delegate, test,
rights and native-environment verification.\citep{exec28_panel_delegates}

\subsection{A common wrapper still carries method-specific choices}

The selected \texttt{nonabsdid} front door dispatches to DCDH, PanelMatch and three \texttt{fect} modes. Preflight verifies
named columns and binary treatment, can encode nonnumeric IDs, rejects entirely missing control columns, and explicitly
warns that partial missingness may yield different effective samples. It does not itself impose a common sample or prove
common estimands. Its PanelMatch path allows reversals and chooses propensity-score matching when controls are supplied,
otherwise no refinement. The \texttt{fect} path uses method-specific tuning defaults: five cross-validation rounds,
an IFE factor cap of two and a matrix-completion lambda grid of five when those values are not supplied. Such changes are
substantive analysis settings and belong in the result provenance.\citep{exec28_panel_delegates}

The DCDH bridge constructs a parsed call with literal horizon arguments and binds data in an evaluation environment. Its
comments attribute this to a particular upstream version's call inspection; that claimed upstream behavior has not been
independently replayed here. A CLI/MCP wrapper should admit known option names and values and pin the upstream interface,
not assume an arbitrary parameter bag is safe or semantically stable. Output conversion also needs care: the effect-cell
helper creates normal intervals from standard errors when bounds are missing, and infers a standard error from interval
width when one is absent. That reverse calculation has the intended meaning for symmetric normal intervals, not arbitrary
bootstrap or simultaneous intervals. The product should retain native bounds and their method, plus label any derived
normal approximation.\citep{exec28_panel_delegates}

\subsection{Stack-count balance and a preserved donor timeline}

The earlier stacked tutorial supplies a useful caution for ports. Its Stata helper sums within-unit row indices. For
$L$ observations of a unit, that sum is $L(L+1)/2$. If every unit in every stack has the same $L$, the factor cancels
from each normalized share: the earlier weights $5/18$ and $15/2$ are unchanged. If two treated units in different stacks
have lengths two and three, while both controls have length two, triangular-count weights are $(2/3,4/3)$; pooled-row
weights are $(4/5,6/5)$; at a common event period with one unit per arm in each stack, weights are $(1,1)$.
These calculations establish that the conventions differ outside balance. They do not establish which unbalanced-panel
estimand is identified. Enforce the declared balanced design or derive and test the alternative explicitly.
\citep{exec28_panel_delegate_reference}

The selected R tutorial's additional comparison to \texttt{TRUE} can include numeric adoption date 1 even if the focal
date is 5 and the clean-control cutoff is 7. A transcribed predicate check admits 1 whereas the intended treated-or-late
rule excludes it. This is a source-level encoding witness, not a native R execution report or an assertion that the
tutorial's calendar-year dataset triggers the branch. Calendar origins and sentinel values need explicit normalization
and rejection tests before the example becomes a general adapter.\citep{exec28_panel_delegate_reference}

For \texttt{pampe}, the selected estimation function fits subset regressions on supplied pre-treatment rows and forms the
counterfactual from the selected OLS model. It rejects missing values and a nonsingle treated-unit request. Its AICc helper
includes the variance parameter in $k$, giving the penalty $2k+2k(k+1)/(n-k-1)$; for $n=8,k=4$ this is $64/3$.
The denominator must remain positive. The top-level path restricts model size, while placebo paths have their own local
choices and require separate checks. Metadata about AICc does not settle all small-sample paths.\citep{exec28_panel_delegates,exec28_panel_delegate_reference}

The leave-one-control-out robustness helper reconstructs pre-period rows as one through the fitted-sample length instead
of retaining the originally supplied indices. A fit on rows 5--12 therefore leads the helper to use rows 1--8. It also refits
the retained donor subset rather than repeating the original full model-selection procedure. These are different
diagnostics. The repair contract is to store exact pre/post row indices and information dates in the result, validate
their reuse, and name fixed-subset sensitivity separately from selection-aware sensitivity. The selected main return
object does not itself retain those original index vectors. An isolated comparison adapter is plausible after these
contracts are repaired and checked; unrestricted adoption or a cross-language port remains premature.
\citep{exec28_panel_delegates,exec28_panel_delegate_reference}

For the bank, these details govern whether a servicing-platform result still refers to the approved cohort, timeline,
target and uncertainty after moving through an agent or another repository. Preserve options, excluded cells, original
IDs, interval provenance and selected model rows. The inspected source permits a concrete integration plan and precise
tests. It does not establish native equivalence, complete dependency rights, empirical coverage or bank acceptance.
